% Options for packages loaded elsewhere
\PassOptionsToPackage{unicode}{hyperref}
\PassOptionsToPackage{hyphens}{url}
%
\documentclass[
  ignorenonframetext,
]{beamer}
\usepackage{pgfpages}
\setbeamertemplate{caption}[numbered]
\setbeamertemplate{caption label separator}{: }
\setbeamercolor{caption name}{fg=normal text.fg}
\beamertemplatenavigationsymbolsempty
% Prevent slide breaks in the middle of a paragraph
\widowpenalties 1 10000
\raggedbottom
\setbeamertemplate{part page}{
  \centering
  \begin{beamercolorbox}[sep=16pt,center]{part title}
    \usebeamerfont{part title}\insertpart\par
  \end{beamercolorbox}
}
\setbeamertemplate{section page}{
  \centering
  \begin{beamercolorbox}[sep=12pt,center]{section title}
    \usebeamerfont{section title}\insertsection\par
  \end{beamercolorbox}
}
\setbeamertemplate{subsection page}{
  \centering
  \begin{beamercolorbox}[sep=8pt,center]{subsection title}
    \usebeamerfont{subsection title}\insertsubsection\par
  \end{beamercolorbox}
}
\AtBeginPart{
  \frame{\partpage}
}
\AtBeginSection{
  \ifbibliography
  \else
    \frame{\sectionpage}
  \fi
}
\AtBeginSubsection{
  \frame{\subsectionpage}
}

\usepackage{amsmath,amssymb}
\usepackage{iftex}
\ifPDFTeX
  \usepackage[T1]{fontenc}
  \usepackage[utf8]{inputenc}
  \usepackage{textcomp} % provide euro and other symbols
\else % if luatex or xetex
  \usepackage{unicode-math}
  \defaultfontfeatures{Scale=MatchLowercase}
  \defaultfontfeatures[\rmfamily]{Ligatures=TeX,Scale=1}
\fi
\usepackage{lmodern}
\ifPDFTeX\else  
    % xetex/luatex font selection
\fi
% Use upquote if available, for straight quotes in verbatim environments
\IfFileExists{upquote.sty}{\usepackage{upquote}}{}
\IfFileExists{microtype.sty}{% use microtype if available
  \usepackage[]{microtype}
  \UseMicrotypeSet[protrusion]{basicmath} % disable protrusion for tt fonts
}{}
\makeatletter
\@ifundefined{KOMAClassName}{% if non-KOMA class
  \IfFileExists{parskip.sty}{%
    \usepackage{parskip}
  }{% else
    \setlength{\parindent}{0pt}
    \setlength{\parskip}{6pt plus 2pt minus 1pt}}
}{% if KOMA class
  \KOMAoptions{parskip=half}}
\makeatother
\usepackage{xcolor}
\newif\ifbibliography
\setlength{\emergencystretch}{3em} % prevent overfull lines
\setcounter{secnumdepth}{-\maxdimen} % remove section numbering


\providecommand{\tightlist}{%
  \setlength{\itemsep}{0pt}\setlength{\parskip}{0pt}}\usepackage{longtable,booktabs,array}
\usepackage{calc} % for calculating minipage widths
\usepackage{caption}
% Make caption package work with longtable
\makeatletter
\def\fnum@table{\tablename~\thetable}
\makeatother
\usepackage{graphicx}
\makeatletter
\newsavebox\pandoc@box
\newcommand*\pandocbounded[1]{% scales image to fit in text height/width
  \sbox\pandoc@box{#1}%
  \Gscale@div\@tempa{\textheight}{\dimexpr\ht\pandoc@box+\dp\pandoc@box\relax}%
  \Gscale@div\@tempb{\linewidth}{\wd\pandoc@box}%
  \ifdim\@tempb\p@<\@tempa\p@\let\@tempa\@tempb\fi% select the smaller of both
  \ifdim\@tempa\p@<\p@\scalebox{\@tempa}{\usebox\pandoc@box}%
  \else\usebox{\pandoc@box}%
  \fi%
}
% Set default figure placement to htbp
\def\fps@figure{htbp}
\makeatother

\makeatletter
\@ifpackageloaded{caption}{}{\usepackage{caption}}
\AtBeginDocument{%
\ifdefined\contentsname
  \renewcommand*\contentsname{Table of contents}
\else
  \newcommand\contentsname{Table of contents}
\fi
\ifdefined\listfigurename
  \renewcommand*\listfigurename{List of Figures}
\else
  \newcommand\listfigurename{List of Figures}
\fi
\ifdefined\listtablename
  \renewcommand*\listtablename{List of Tables}
\else
  \newcommand\listtablename{List of Tables}
\fi
\ifdefined\figurename
  \renewcommand*\figurename{Figure}
\else
  \newcommand\figurename{Figure}
\fi
\ifdefined\tablename
  \renewcommand*\tablename{Table}
\else
  \newcommand\tablename{Table}
\fi
}
\@ifpackageloaded{float}{}{\usepackage{float}}
\floatstyle{ruled}
\@ifundefined{c@chapter}{\newfloat{codelisting}{h}{lop}}{\newfloat{codelisting}{h}{lop}[chapter]}
\floatname{codelisting}{Listing}
\newcommand*\listoflistings{\listof{codelisting}{List of Listings}}
\makeatother
\makeatletter
\makeatother
\makeatletter
\@ifpackageloaded{caption}{}{\usepackage{caption}}
\@ifpackageloaded{subcaption}{}{\usepackage{subcaption}}
\makeatother

\usepackage{bookmark}

\IfFileExists{xurl.sty}{\usepackage{xurl}}{} % add URL line breaks if available
\urlstyle{same} % disable monospaced font for URLs
\hypersetup{
  pdftitle={Incorporating AI  into Teaching},
  pdfauthor={Kerry Back Rice University FMA 2025},
  hidelinks,
  pdfcreator={LaTeX via pandoc}}


\title{Incorporating AI into Teaching}
\author{Kerry BackRice UniversityFMA 2025}
\date{}

\begin{document}
\frame{\titlepage}


\begin{frame}{What the Future Holds (CNBC, 9/30/2025)}
\phantomsection\label{what-the-future-holds-cnbc-9302025}
Derek Waldron, JPMorgan chief analytics officer:

The broad vision that we're working towards is one where the JPMorgan
Chase of the future is going to be a fully AI-connected enterprise.

Every employee will have their own personalized AI assistant; every
process is powered by AI agents, and every client experience has an AI
concierge.
\end{frame}

\begin{frame}{}
\phantomsection\label{section}
Waldron gave {[}JPMorgan LLM Suite{]} a prompt:

You are a technology banker at JPMorgan Chase preparing for a meeting
with the CEO and CFO of Nvidia. Prepare a five-page presentation that
includes the latest news, earnings and a peer comparison.

LLM Suite created a credible-looking PowerPoint deck in about 30
seconds.
\end{frame}

\begin{frame}{}
\phantomsection\label{section-1}
Waldron:

You can imagine in the past how that would have been done; we would've
had teams of investment banking analysts working long hours at night to
do this.

\begin{block}{From chat to agents}
\phantomsection\label{from-chat-to-agents}
JPMorgan is now early in the next phase of its AI blueprint: It has
begun deploying agentic AI to handle complex multistep tasks for
employees, according to an internal road map provided by the bank.
\end{block}
\end{frame}

\begin{frame}{Life was the same for 40 years}
\phantomsection\label{life-was-the-same-for-40-years}
\begin{columns}[T]
\begin{column}{0.33\linewidth}
\includegraphics[width=3.125in,height=2.08333in]{images/visicalc.png}

VisiCalc introduced for the Apple II in 1979
\end{column}

\begin{column}{0.33\linewidth}
\includegraphics[width=3.125in,height=2.08333in]{images/lotus-123.jpg}

Lotus-123 introduced for the IBM PC in 1983
\end{column}

\begin{column}{0.33\linewidth}
\includegraphics[width=3.125in,height=2.08333in]{images/excel20.png}

Microsoft Excel introduced for Windows in 1987
\end{column}
\end{columns}
\end{frame}

\begin{frame}{Now, with AI, we're all in the same boat}
\phantomsection\label{now-with-ai-were-all-in-the-same-boat}
\pandocbounded{\includegraphics[keepaspectratio]{images/ducks-sinking.png}}
\end{frame}

\begin{frame}{It's hard to keep up!}
\phantomsection\label{its-hard-to-keep-up}
\pandocbounded{\includegraphics[keepaspectratio]{images/coyote11.png}}

\pandocbounded{\includegraphics[keepaspectratio]{images/coyote21.png}}
\end{frame}

\begin{frame}{Plan}
\phantomsection\label{plan}
\begin{enumerate}
\tightlist
\item
  Example of what AI can do today - new Claude skill from September
\item
  Historical analogues?
\item
  A course on AI
\item
  Updating existing courses to use AI
\item
  Deep dive into Claude Code
\end{enumerate}
\end{frame}

\begin{frame}{1. Example (prompt for Claude)}
\phantomsection\label{example-prompt-for-claude}
Read the uploaded case (Valuing Walmart 2010). Ignore the valuation
method described in the case. Instead, generate pro forma statements and
perform a two-stage DCF analysis.

Document your assumptions and your reasons for them. Generate an Excel
workbook and a Word doc. Format the Excel worksheets nicely.

\emph{Currently requires Claude Max subscription.}
\end{frame}

\begin{frame}{Microsoft Vibe Working (powered by Anthropic)}
\phantomsection\label{microsoft-vibe-working-powered-by-anthropic}
\end{frame}

\begin{frame}{2. Historical analogues}
\phantomsection\label{historical-analogues}
\begin{enumerate}
\tightlist
\item
  Handheld calculators
\item
  Personal computers
\end{enumerate}

How did teaching and learning change? Are there lessons for today?
\end{frame}

\begin{frame}{Introduction of Handheld Calculators}
\phantomsection\label{introduction-of-handheld-calculators}
\begin{itemize}
\tightlist
\item
  Much worry that students wouldn't learn addition, subtraction,
  \ldots{}

  \begin{itemize}
  \tightlist
  \item
    Studies haven't found an effect
  \item
    But primary schools still provide practice and assessments in which
    calculators are not allowed
  \end{itemize}
\item
  Some hoped that schools could move beyond calculation and teach deeper
  math

  \begin{itemize}
  \tightlist
  \item
    That didn't seem to happen either
  \end{itemize}
\end{itemize}
\end{frame}

\begin{frame}{Introduction of Personal computers \& spreadsheets}
\phantomsection\label{introduction-of-personal-computers-spreadsheets}
\begin{enumerate}
\tightlist
\item
  We added courses on how to use spreadsheets
\item
  We updated existing courses to use spreadsheets
\item
  We assess how well students can use spreadsheets (no pencil and paper
  valuation analyses)
\end{enumerate}

\begin{quote}
Teaching how to do finance in spreadsheets is a main goal.
\end{quote}
\end{frame}

\begin{frame}{Why?}
\phantomsection\label{why}
\begin{itemize}
\tightlist
\item
  Spreadsheets are used in business
\item
  Spreadsheets are ideal for teaching the logic of some things
  (valuation/capital budgeting)
\item
  Spreadsheets don't automatically do the things we were teaching before
  (financial calculators \(\rightarrow\) spreadsheets)
\end{itemize}
\end{frame}

\begin{frame}{What about AI?}
\phantomsection\label{what-about-ai}
✓ AI is used in business

✓ AI is useful for teaching/tutoring

✗ AI can automatically do many things we were teaching before

AI is more like calculators than spreadsheets. We worry that students
won't learn skills that they need.

Like primary schools and arithmetic, we will need to have practice and
assessments in which AI is not allowed
\end{frame}

\begin{frame}{AI is used in Business}
\phantomsection\label{ai-is-used-in-business}
MIT/Media Labs State of AI in Business 2025:

Workers from over 90\% of the companies we surveyed reported regular use
of personal AI tools for work tasks. In fact, almost every single person
used an LLM in some form for their work.
\end{frame}

\begin{frame}{Plan}
\phantomsection\label{plan-1}
✓ Example of what AI can do today

✓ Historical analogues

\begin{enumerate}
\setcounter{enumi}{2}
\item
  A course on AI
\item
  Updating existing courses to use AI
\item
  Deep dive into Claude Code
\end{enumerate}
\end{frame}

\begin{frame}{3. Course on AI: principles}
\phantomsection\label{course-on-ai-principles}
\begin{itemize}
\tightlist
\item
  {Treat AI as a colleague:} Collaborate, evaluate, iterate
\item
  {Vibe code financial analyses: LLMs are not trustworthy for
  arithmetic:} Coding (Python) is essential
\item
  {Create apps:} For critical operations, test code and save as app
\item
  Create custom chatbots:\textless{}\span\textgreater{} Chatbots can be
  customized through the system prompt, RAG, or fine tuning
\item
  {Create agents:} To use chatbot for critical operations, create app
  and provide as tool
\end{itemize}
\end{frame}

\begin{frame}{Prompt Engineering}
\phantomsection\label{prompt-engineering}
\begin{itemize}
\tightlist
\item
  Some students ask for instruction on how to prompt AI
\item
  Special techniques aren't needed. Practice is.
\end{itemize}

Ethan Mollick, Wharton, ``Good Enough Prompting:''

Treat AI like an infinitely patient new coworker \ldots{} As it is a
coworker, you want to work with it, not just give it orders, and you
also want to learn what it is good or bad at \ldots{} Working with AI is
a dialogue, not an order.
\end{frame}

\begin{frame}{Course structure}
\phantomsection\label{course-structure}
A. Introduction

B. Coding with AI

C. Apps, Chatbots, and Agents
\end{frame}

\begin{frame}{A. Course Introduction}
\phantomsection\label{a.-course-introduction}
\begin{itemize}
\tightlist
\item
  Example of working with AI on valuation: Claude/Excel example
\item
  Case on corporate implementation of AI
\end{itemize}
\end{frame}

\begin{frame}{Example of working with AI}
\phantomsection\label{example-of-working-with-ai}
Claude/Excel example

\begin{itemize}
\tightlist
\item
  Explain the calculations
\item
  What did you assume for \ldots{} and why?
\item
  How does the value change if we assume \ldots{} instead?
\end{itemize}
\end{frame}

\begin{frame}{Case on corporate implementation}
\phantomsection\label{case-on-corporate-implementation}
\href{https://store.hbr.org/product/blazing-new-trails-responsible-generative-ai-and-the-creative-adoption-of-a-large-language-model-at-deloitte-canada/HEC382?srsltid=AfmBOoo8whYuCV6aPeMolvqA2FNyJJz0-l8D3oFSoKKObmqppQ_Y_QcO}{Blazing
New Trails: Responsible Generative AI and the Creative Adoption of a
Large Language Model at Deloitte Canada}, 2024.

\begin{itemize}
\tightlist
\item
  Describes creation of a custom chatbot
\item
  Discusses reliability, bias, compliance, data privacy, trust of
  customers
\end{itemize}
\end{frame}

\begin{frame}{B. Coding with AI}
\phantomsection\label{b.-coding-with-ai}
In March 2025, Y Combinator CEO Garry Tan and managing partner Jared
Friedman stated that

\begin{columns}[T]
\begin{column}{0.48\linewidth}
for roughly a quarter of the startups in their Winter 2025 cohort, 95\%
of the codebase was written by AI.
\end{column}

\begin{column}{0.48\linewidth}
\pandocbounded{\includegraphics[keepaspectratio]{images/bugs-bunny.png}}
\end{column}
\end{columns}
\end{frame}

\begin{frame}{Vibe coding for finance}
\phantomsection\label{vibe-coding-for-finance}
\begin{itemize}
\tightlist
\item
  Chatbots with python can easily do many types of financial analyses:
  regression, mean-variance optimization, binomial-tree option valuation
  \ldots{}
\item
  Running python in the cloud:

  \begin{itemize}
  \tightlist
  \item
    ChatGPT, Claude, Google Colab with Gemini, Julius.ai
  \end{itemize}
\item
  Running python locally:

  \begin{itemize}
  \tightlist
  \item
    VS Code, Cursor, Windsurf, Claude Code
  \end{itemize}
\end{itemize}
\end{frame}

\begin{frame}{C. Apps, Chatbots, and Agents}
\phantomsection\label{c.-apps-chatbots-and-agents}
\begin{columns}[T]
\begin{column}{0.7\linewidth}
\pandocbounded{\includegraphics[keepaspectratio]{images/ai_learning_progression.png}}
\end{column}

\begin{column}{0.3\linewidth}
\begin{itemize}
\tightlist
\item
  Apps can use AI
\item
  Chatbots are apps
\item
  Agents are chatbots with tools
\end{itemize}
\end{column}
\end{columns}
\end{frame}

\begin{frame}{Building apps, chatbots, and agents}
\phantomsection\label{building-apps-chatbots-and-agents}
\begin{itemize}
\tightlist
\item
  Point and click to build Custom GPTs at chat.openai.com
\item
  Vibe code
\end{itemize}
\end{frame}

\begin{frame}{chatbot is an app consisting of}
\phantomsection\label{chatbot-is-an-app-consisting-of}
\begin{itemize}
\tightlist
\item
  User interface
\item
  API connection to an LLM, which has possibly been fine-tuned
\item
  System prompt
\item
  Possible retrieval of documents
\item
  Possible use of tools
\end{itemize}
\end{frame}

\begin{frame}{System prompt is fundamental}
\phantomsection\label{system-prompt-is-fundamental}
A system prompt is text that is sent to the LLM along with each user
prompt.

It contains information and instructions for the LLM.

\href{https://docs.claude.com/en/release-notes/system-prompts\#august-5-2025}{Anthropic's
System Prompts}
\end{frame}

\begin{frame}{RAG (Retrieval Augmented Generation) chatbots}
\phantomsection\label{rag-retrieval-augmented-generation-chatbots}
\begin{itemize}
\tightlist
\item
  Send stored chunks of documents to LLM along with prompt and system
  prompt
\item
  Choose chunks most relevant to the prompt

  \begin{itemize}
  \tightlist
  \item
    Vector embedding (vector store)
  \item
    Similarity search
  \end{itemize}
\item
  Key element of many corporate AI implementations
\end{itemize}
\end{frame}

\begin{frame}{Chatbot structure}
\phantomsection\label{chatbot-structure}
\pandocbounded{\includegraphics[keepaspectratio]{images/chatbot.png}}
\end{frame}

\begin{frame}{Agent = Chatbot + Tools}
\phantomsection\label{agent-chatbot-tools}
\begin{itemize}
\tightlist
\item
  AI agent is a chatbot equipped with tools to do things
\item
  Claude and ChatGPT are agents - they can send code to Python
  environments
\item
  Database tools are very useful for chatbots
\end{itemize}
\end{frame}

\begin{frame}{Claude for Financial Services}
\phantomsection\label{claude-for-financial-services}
From
\href{https://www.anthropic.com/news/claude-for-financial-services}{Anthropic}:

Through data providers, Claude has real-time access to comprehensive
financial information:

\begin{columns}[T]
\begin{column}{0.33\linewidth}
\begin{itemize}
\tightlist
\item
  Box
\item
  Daloopa
\item
  Databricks
\end{itemize}
\end{column}

\begin{column}{0.33\linewidth}
\begin{itemize}
\tightlist
\item
  FactSet
\item
  Morningstar
\item
  Palantir
\end{itemize}
\end{column}

\begin{column}{0.33\linewidth}
\begin{itemize}
\tightlist
\item
  PitchBook
\item
  S\&P Global
\item
  Snowflake
\end{itemize}
\end{column}
\end{columns}
\end{frame}

\begin{frame}{Why mba students should learn about agents}
\phantomsection\label{why-mba-students-should-learn-about-agents}
Fortune, 9-14-2025:

PromptQL, an enterprise AI platform created by San Francisco-based
developer tooling company Hasura, is doling out {\$900-per-hour} wages
to its engineers tasked with building and deploying AI agents to analyze
internal company data using large language models (LLMs).

Tanmai Gopal, PromptQL's cofounder and CEO, said ``{MBA types} \ldots{}
are very strategic thinkers, and they're smart people, but they {don't
have an intuition for what AI can do}.''
\end{frame}

\begin{frame}{database agent passes prompt from user to LLM}
\phantomsection\label{database-agent-passes-prompt-from-user-to-llm}
\begin{columns}[T]
\begin{column}{0.6\linewidth}
\pandocbounded{\includegraphics[keepaspectratio]{images/agent_v1.png}}
\end{column}

\begin{column}{0.4\linewidth}
May require chatting at this stage to clarify user's request
\end{column}
\end{columns}
\end{frame}

\begin{frame}{agent passes SQL from LLM to data warehouse}
\phantomsection\label{agent-passes-sql-from-llm-to-data-warehouse}
\begin{columns}[T]
\begin{column}{0.6\linewidth}
\pandocbounded{\includegraphics[keepaspectratio]{images/agent_v2.png}}
\end{column}

\begin{column}{0.4\linewidth}
\end{column}
\end{columns}
\end{frame}

\begin{frame}{Data Warehouse passes response to agent}
\phantomsection\label{data-warehouse-passes-response-to-agent}
\begin{columns}[T]
\begin{column}{0.6\linewidth}
\pandocbounded{\includegraphics[keepaspectratio]{images/agent_v3.png}}
\end{column}

\begin{column}{0.4\linewidth}
\begin{itemize}
\tightlist
\item
  Response could be error message
\item
  If so, agent should send to LLM for new SQL code
\end{itemize}
\end{column}
\end{columns}
\end{frame}

\begin{frame}{agent passes data from warehouse to user}
\phantomsection\label{agent-passes-data-from-warehouse-to-user}
\begin{columns}[T]
\begin{column}{0.6\linewidth}
\pandocbounded{\includegraphics[keepaspectratio]{images/agent_v4.png}}
\end{column}

\begin{column}{0.4\linewidth}
\begin{itemize}
\tightlist
\item
  Could be more things connected
\item
  E.g., send data to python engine
\item
  Send python output to user
\end{itemize}
\end{column}
\end{columns}
\end{frame}

\begin{frame}{What can we do in the classroom?}
\phantomsection\label{what-can-we-do-in-the-classroom}
\begin{itemize}
\tightlist
\item
  We can show students how to vibe code apps to automate things

  \begin{itemize}
  \tightlist
  \item
    Getting data
  \item
    Financial analysis
  \item
    Getting written analysis from AI
  \item
    Producing Excel, Word, PowerPoint
  \end{itemize}
\item
  And convert app to chatbot tool
\end{itemize}
\end{frame}

\begin{frame}{Tools}
\phantomsection\label{tools}
\begin{itemize}
\tightlist
\item
  Point and click: ChatGPT's Custom GPTs
\item
  Simplest vibe coding environment: Replit
\item
  Julius.ai can build but not run Streamlit apps
\item
  Google Colab can build Streamlit apps and deploy via ngrok
\item
  Permanent deployment: Koyeb, Render, \ldots{}
\end{itemize}

How do we learn this stuff to teach it: Ask AI
\end{frame}

\begin{frame}{Creating openai custom gpt (No Code)}
\phantomsection\label{creating-openai-custom-gpt-no-code}
\begin{itemize}
\tightlist
\item
  \textbf{Configure}: Name, description, and system prompt
\item
  \textbf{Add capabilities}: Web browsing, DALL-E, Code Interpreter
\item
  \textbf{Upload knowledge}: Documents for RAG (PDFs, text files)
\item
  \textbf{Add actions}: Connect to external APIs (optional)
\item
  \textbf{Test \& publish}: Private, link-only, or public to GPT Store
\end{itemize}

\emph{Requires ChatGPT Plus subscription}
\end{frame}

\begin{frame}{Plan}
\phantomsection\label{plan-2}
✓ Example of what AI can do today

✓ Historical analogues

✓ A course on AI

\begin{enumerate}
\setcounter{enumi}{3}
\item
  Updating existing courses to use AI
\item
  Deep dive into Claude Code
\end{enumerate}
\end{frame}

\begin{frame}{4. Updating existing courses: what we should teach}
\phantomsection\label{updating-existing-courses-what-we-should-teach}
\begin{itemize}
\tightlist
\item
  We generally teach concepts first with slides (or board)
\item
  Then show in tools (usually Excel)
\item
  Now, there is a next step: chatbots and agents
\end{itemize}

Teach about implementation in Excel \(\rightarrow\) students build Excel
models.

Teach about implementation in AI \(\rightarrow\) students build prompts
or apps
\end{frame}

\begin{frame}{How we should teach}
\phantomsection\label{how-we-should-teach}
\begin{itemize}
\tightlist
\item
  AI is a very effective tutor
\item
  We should take advantage of it
\item
  {Asking students to chat with chatbots can replace ``Do the assigned
  reading.''}
\item
  Can be assigned before class or in class before slides.
\end{itemize}

Ask your chatbot to teach you about \ldots{} Be sure to tell it to ask
you questions to ensure you are understanding.
\end{frame}

\begin{frame}{Controlling the sources}
\phantomsection\label{controlling-the-sources}
\begin{itemize}
\tightlist
\item
  Build RAG (Google's Notebook LM or Custom GPT or \ldots)
\item
  If sources are copyrighted, first ask AI for a complete report on the
  subject, read and edit the AI report, and add to system prompt ``Teach
  from this report \ldots{}''
\end{itemize}

But this is probably unnecessary for most topics. The LLMs are much,
much more reliable than they were a couple of years ago when everyone
was concerned about hallucinations.
\end{frame}

\begin{frame}{How can we assess?}
\phantomsection\label{how-can-we-assess}
\begin{itemize}[<+->]
\tightlist
\item
  Out-of-class assignments should become ``turn in your chat, its
  output, and your evaluation of the output'' or ``turn in your app.''
\item
  In-class {exams without AI} to ensure students understand concepts
\item
  In-class {exams with AI} that are more ambitious than exams today.
  Cases can become exams with AI assistance.
\end{itemize}
\end{frame}

\begin{frame}{Plan}
\phantomsection\label{plan-3}
✓ Example of what AI can do today

✓ Historical analogues

✓ A course on AI

✓ Updating existing courses to use AI

\begin{enumerate}
\setcounter{enumi}{4}
\tightlist
\item
  Deep dive into Claude Code
\end{enumerate}
\end{frame}

\begin{frame}{Claude Code}
\phantomsection\label{claude-code}
\begin{itemize}
\tightlist
\item
  Terminal based coding agent - also OpenAI Codex, Gemini CLI
\item
  Can be used as a plug-in to VS Code, Cursor, Windsurf
\item
  Uses standard Anthropic models, but seems to have a great system
  prompt
\end{itemize}

\emph{Requires Claude Pro subscription.}
\end{frame}

\begin{frame}{Capabilities}
\phantomsection\label{capabilities}
\begin{itemize}
\tightlist
\item
  Create, edit, copy, move, delete files
\item
  Write, edit, and run Python scripts and notebooks
\item
  Write, edit, and compile LaTeX documents
\item
  Create and edit Word docs and PowerPoint decks
\item
  Create GitHub repos and push and launch apps on hosts (Koyeb, Render,
  \ldots)
\end{itemize}
\end{frame}

\begin{frame}{Claude Code Demo}
\phantomsection\label{claude-code-demo}
\begin{itemize}
\tightlist
\item
  Create a streamlit app for mean-variance optimization. The user will
  input the number of assets, their means, standard deviations, and
  correlations and risk-free rate.
\item
  Compute and display the tangency portfolio. Run the app.
\end{itemize}

Then: show the code that calculates the tangency portfolio and explain
it.
\end{frame}

\begin{frame}{Claude Code in VS Code, etc.}
\phantomsection\label{claude-code-in-vs-code-etc.}
\begin{itemize}
\tightlist
\item
  VS Code, Cursor, Windsurf are full-featured Python and LaTeX IDEs
\item
  Examples of Claude Code for LaTeX and Python:

  \begin{itemize}
  \tightlist
  \item
    Read this paragraph and suggest improvements .
  \item
    Make these global edits \ldots{}
  \item
    Do this analysis and generate a figure or LaTeX table and insert it
    into the document \ldots{}
  \end{itemize}
\end{itemize}
\end{frame}

\begin{frame}{Aaron Linsky, CTO, AIA Labs at Bridgewater:}
\phantomsection\label{aaron-linsky-cto-aia-labs-at-bridgewater}
We've been developing capabilities powered by Claude since 2023 within
AIA Labs.

Claude powered the first versions of our {Investment Analyst Assistant},
which streamlined our analysts' workflow by generating Python code,
creating data visualizations, and iterating through { complex financial
analysis tasks with the precision of a junior analyst}.
\end{frame}

\begin{frame}{Yeyati, Brookings, 2025:}
\phantomsection\label{yeyati-brookings-2025}
As AI models begin to handle underwriting, compliance, and asset
allocation, the traditional architecture of financial work is undergoing
a fundamental shift.

As job descriptions evolve, so does the definition of financial talent.
{Excel is no longer a differentiator. Python is fast becoming the new
Excel. }

But technical skills alone will not cut it. The {most in demand profiles
today are those that speak both AI and finance}.
\end{frame}

\begin{frame}{}
\phantomsection\label{section-2}
\end{frame}

\begin{frame}{}
\phantomsection\label{final-slide}
\pandocbounded{\includegraphics[keepaspectratio]{images/qr_code_slides16x9.png}}

\href{https://fma2025.kerryback.com}{fma2025.kerryback.com}
\end{frame}

\begin{frame}{Example Data Chatbot (fully vibe coded)}
\phantomsection\label{example-data-chatbot-fully-vibe-coded}
\end{frame}




\end{document}
